% GENERATED FILE. Edit canonical lessons, then run npm run generate:books.
% canonical-source-hash: fnv1a64:91f3240b2b29f732
% canonical-lessons: FA-C116-keshavarz, FA-C116-ranande, FA-C116-furushande, FA-C116-nanva, FA-C116-kudak

\chapter{Farmer, Driver, Shopkeeper, Baker, Small child}
\label{ch:fa-farmer-driver-shopkeeper-baker-small-child}
\hlchaptermodality{116}

\begin{chapteropening}
\textbf{By the end of this chapter:} \emph{I can say a farmer, a driver, a shopkeeper, a baker and a small child.}

Complete the last lesson of chapter 116: I can say a farmer, a driver, a shopkeeper, a baker and a small child.
\end{chapteropening}

\section[\texorpdfstring{\fa{کشاورز} (keshâvarz)}{کشاورز (keshâvarz)}]{\fa{کشاورز} (keshâvarz) — a farmer}

\label{lesson:FA-C116-keshavarz}



\begin{quote}
Before the new one: say the Persian for a cook, then the Persian for a worker.
\end{quote}

\subsection*{You'll want to know: \fa{کشاورز}}
\textbf{\fa{کشاورز}} — \emph{keshâvarz} — \textquotedblleft{}a farmer\textquotedblright{}.

\begin{cousinweb}
\emph{keshâvarz} is a native Persian word, inherited rather than borrowed; it is usually traced to Middle Persian \emph{kišāwarz}, \textquotedblleft{}farmer\textquotedblright{}. It is usually explained as built on \emph{keš-}, \textquotedblleft{}to plough\textquotedblright{}.
\end{cousinweb}

\begin{grammarlens}[title={What you've built}]
One more word for this chapter.
\end{grammarlens}

\subsection*{Your turn}
\begin{itemize}
\raggedright
  \item \emph{Say it:} \emph{keshâvarz}
  \item \emph{Say it:} \emph{keshâvarz}, once more
  \item \emph{Say it:} say \emph{keshâvarz}
  \item \emph{Recall:} say the Persian for a cook, then the Persian for a worker, then say \emph{keshâvarz} again
\end{itemize}

\begin{tcolorbox}[breakable,colback=teal!4,colframe=teal!35!black,title={Before you move on}]
What does \fa{کشاورز} mean? (\textquotedblleft{}A farmer\textquotedblright{}.) Say it once more.
\end{tcolorbox}
\section[\texorpdfstring{\fa{راننده} (rânande)}{راننده (rânande)}]{\fa{راننده} (rânande) — a driver}

\label{lesson:FA-C116-ranande}



\begin{quote}
Before the new one: say the Persian for a worker, then the Persian for a farmer.
\end{quote}

\subsection*{You'll want to know: \fa{راننده}}
\textbf{\fa{راننده}} — \emph{rânande} — \textquotedblleft{}a driver\textquotedblright{}.

\begin{cousinweb}
\emph{rânande} is the \textquotedblleft{}-ing\textquotedblright{} form of \emph{rândan}, \textquotedblleft{}to drive\textquotedblright{}: \textquotedblleft{}the one driving\textquotedblright{}. The verb is inherited, usually traced to Middle Persian \emph{rāndan}.
\end{cousinweb}

\begin{grammarlens}[title={What you've built}]
One more word for this chapter.
\end{grammarlens}

\subsection*{Your turn}
\begin{itemize}
\raggedright
  \item \emph{Say it:} \emph{rânande}
  \item \emph{Say it:} \emph{rânande}, once more
  \item \emph{Say it:} say \emph{rânande}
  \item \emph{Recall:} say the Persian for a worker, then the Persian for a farmer, then say \emph{rânande} again
\end{itemize}

\begin{tcolorbox}[breakable,colback=teal!4,colframe=teal!35!black,title={Before you move on}]
What does \fa{راننده} mean? (\textquotedblleft{}A driver\textquotedblright{}.) Say it once more.
\end{tcolorbox}
\section[\texorpdfstring{\fa{فروشنده} (furushande)}{فروشنده (furushande)}]{\fa{فروشنده} (furushande) — a shopkeeper}

\label{lesson:FA-C116-furushande}



\begin{quote}
Before the new one: say the Persian for a farmer, then the Persian for a driver.
\end{quote}

\subsection*{You'll want to know: \fa{فروشنده}}
\textbf{\fa{فروشنده}} — \emph{furushande} — \textquotedblleft{}a shopkeeper\textquotedblright{}.

\begin{cousinweb}
\emph{furushande} is \emph{furush-}, the present stem of \emph{forukhtan}, \textquotedblleft{}to sell\textquotedblright{}, with \emph{-ande}, \textquotedblleft{}-ing\textquotedblright{}: \textquotedblleft{}the one selling\textquotedblright{}. The verb is inherited, usually traced to Middle Persian \emph{frōxtan}.
\end{cousinweb}

\begin{grammarlens}[title={What you've built}]
One more word for this chapter.
\end{grammarlens}

\subsection*{Your turn}
\begin{itemize}
\raggedright
  \item \emph{Say it:} \emph{furushande}
  \item \emph{Say it:} \emph{furushande}, once more
  \item \emph{Say it:} say \emph{furushande}
  \item \emph{Recall:} say the Persian for a farmer, then the Persian for a driver, then say \emph{furushande} again
\end{itemize}

\begin{tcolorbox}[breakable,colback=teal!4,colframe=teal!35!black,title={Before you move on}]
What does \fa{فروشنده} mean? (\textquotedblleft{}A shopkeeper\textquotedblright{}.) Say it once more.
\end{tcolorbox}
\section[\texorpdfstring{\fa{نانوا} (nânvâ)}{نانوا (nânvâ)}]{\fa{نانوا} (nânvâ) — a baker}

\label{lesson:FA-C116-nanva}



\begin{quote}
Before the new one: say the Persian for a driver, then the Persian for a shopkeeper.
\end{quote}

\subsection*{You'll want to know: \fa{نانوا}}
\textbf{\fa{نانوا}} — \emph{nânvâ} — \textquotedblleft{}a baker\textquotedblright{}.

\begin{cousinweb}
\emph{nânvâ} is built on \emph{nân}, \textquotedblleft{}bread\textquotedblright{}, an inherited word: EtymDB traces it to Middle Persian \emph{nān}. The ending \emph{-vâ} marks the one who deals in it.
\end{cousinweb}

\begin{grammarlens}[title={What you've built}]
One more word for this chapter.
\end{grammarlens}

\subsection*{Your turn}
\begin{itemize}
\raggedright
  \item \emph{Say it:} \emph{nânvâ}
  \item \emph{Say it:} \emph{nânvâ}, once more
  \item \emph{Say it:} say \emph{nânvâ}
  \item \emph{Recall:} say the Persian for a driver, then the Persian for a shopkeeper, then say \emph{nânvâ} again
\end{itemize}

\begin{tcolorbox}[breakable,colback=teal!4,colframe=teal!35!black,title={Before you move on}]
What does \fa{نانوا} mean? (\textquotedblleft{}A baker\textquotedblright{}.) Say it once more.
\end{tcolorbox}
\section[\texorpdfstring{\fa{کودک} (kudak)}{کودک (kudak)}]{\fa{کودک} (kudak) — a small child}

\label{lesson:FA-C116-kudak}



\begin{quote}
Before the new one: say the Persian for a shopkeeper, then the Persian for a baker.
\end{quote}

\subsection*{You'll want to know: \fa{کودک}}
\textbf{\fa{کودک}} — \emph{kudak} — \textquotedblleft{}a small child\textquotedblright{}.

\begin{cousinweb}
\emph{kudak} is a native Persian word, inherited rather than borrowed; it is usually traced to Middle Persian \emph{kōdak}, \textquotedblleft{}small child\textquotedblright{}.
\end{cousinweb}

\begin{grammarlens}[title={What you've built}]
One more word for this chapter.
\end{grammarlens}

\subsection*{Your turn}
\begin{itemize}
\raggedright
  \item \emph{Say it:} \emph{kudak}
  \item \emph{Say it:} \emph{kudak}, once more
  \item \emph{Say it:} say \emph{kudak}
  \item \emph{Recall:} say the Persian for a shopkeeper, then the Persian for a baker, then say \emph{kudak} again
\end{itemize}

\begin{tcolorbox}[breakable,colback=teal!4,colframe=teal!35!black,title={Before you move on}]
What does \fa{کودک} mean? (\textquotedblleft{}A small child\textquotedblright{}.) Say it once more.
\end{tcolorbox}
